\section*{第 \ref{ch:b2:guided-waves} 章——导行波与谐振腔}

\begin{solution}{exo:b2:guided-waves:1}
$f_n = nc/2a = 7.5$、$15$、\qty{22.5}{GHz}。在 \qty{12}{GHz} 处只有 $n = 1$：$k_g =
(2\pi/c)\sqrt{f^2 - f_1^2} = \qty{196}{rad/m}$，$\lambda_g = \qty{3.2}{cm}$，$v_\varphi = c/0.78
= \qty{3.8e8}{m/s}$，$v_g = \qty{2.3e8}{m/s}$。在 \qty{6}{GHz} 处：$\ell = c/2\pi\sqrt{f_1^2 -
f^2} = \qty{1.1}{cm}$。
\end{solution}

\begin{solution}{exo:b2:guided-waves:2}
TE$_{10}$ \qty{6.55}{GHz}，TE$_{20}$ \qty{13.1}{GHz}，TE$_{01}$ \qty{14.7}{GHz}：单模工作
区为 $6.55$ 至 \qty{13.1}{GHz}。在 \qty{10}{GHz} 处：$\sqrt{1 - (6.55/10)^2} =
0.756$，$\lambda_g = \qty{4.0}{cm}$，$v_g = 0.76c$，$\cos\theta = 0.756$，$\theta = \qty{41}{\degree}$。
$b \approx a/2$ 使 TE$_{01}$ 仍高于 TE$_{20}$（最宽的单模频带），
同时为场留出最大间隙（决定击穿前的功率容量）。
\end{solution}

\begin{solution}{exo:b2:guided-waves:3}
$f = 1.5 \times 10^8\sqrt{(m/0.3)^2 + (n/0.3)^2 + (p/0.2)^2}$：\qty{0.71}{GHz}、\qty{0.90}{GHz}、
\qty{1.35}{GHz}、\qty{1.95}{GHz}。$N = 8\pi Vf^3/3c^3 = 8\pi \times 0.018 \times 1.56 \times 10^{28}/
8.1 \times 10^{25} \approx 260$。立方腔：$f = c\sqrt2/2L$，$L = \qty{2.3}{cm}$。
\end{solution}

\begin{solution}{exo:b2:guided-waves:4}
$\mathrm{NA} = \sqrt{1.48^2 - 1.46^2} = 0.24$，\qty{14}{\degree}；$\Delta t = (n_1L/c)(n_1/n_2
- 1) = \qty{68}{ns/km}$；\qty{2}{km} 上为 \qty{135}{ns}：\qty{7}{Mbit/s}；取 $0.003$ 时：
\qty{10}{ns/km}，\qty{50}{Mbit/s}。
\end{solution}

\begin{solution}{exo:b2:guided-waves:5}
(a) $\underline B_y = (k_g/\omega)\underline E$，$\underline B_z = (\iu/\omega)\partial_y\underline E = (\iu\pi/a\omega)
E_0\cos(\pi y/a)\eu^{\iu(\omega t - k_gz)}$：二者相位正交。(b) 在 $y = 0$ 处，$B_y = 0$，且
$\vect j_s = \vect e_y\wedge\vect B/\mu_0 = (B_z/\mu_0)\vect e_x$：沿场的方向。(c)
$\langle\Pi_z\rangle = E_0^2k_g\sin^2(\pi y/a)/2\mu_0\omega$，单位宽度的功率
为 $E_0^2k_ga/4\mu_0\omega$；$\langle\Pi_y\rangle \propto \operatorname{Re}(\underline E\,\iu\underline B_z^*)
= 0$。(d) 单位长度与单位宽度上的平均能量为 $\tfrac14\varepsilon_0E_0^2a$（电
场与磁场各占 $\tfrac18\varepsilon_0E_0^2a$）；二者之比为 $c^2k_g/\omega = v_g$。
\end{solution}

\begin{solution}{exo:b2:guided-waves:6}
(a) $f_1 = \qty{150}{GHz} \gg f$：$\ell \approx a/\pi = \qty{0.32}{mm}$；$\eu^{-0.5/0.32}$：
按振幅计为 \qty{-14}{dB}，按功率计为 $-27$。(b) $f_1 = \qty{750}{MHz}$；在 \qty{100}{MHz} 处
$\ell = \qty{6.4}{cm}$ 时，一米内 \qty{-136}{dB}。(c) $f_1 = \qty{15}{GHz}$，$\ell =
\qty{3.2}{mm}$：\qty{27}{dB/cm}。(d) 波无法传播；
倏逝性源于几何结构，而非耗散。
\end{solution}

\begin{solution}{exo:b2:guided-waves:7}
(a) $P = \tfrac14\varepsilon_0cE_0^2ab\sqrt{1 - f_{10}^2/f^2} = \qty{120}{kW}$。(b) 放大 $\times9$：
\qty{1}{MW}；加压氮气或 SF$_6$ 可提高击穿场强。(c)
$j_s \approx E_0/\mu_0c = \qty{2.7}{kA/m}$；$\tfrac12R_sj_s^2 = \qty{9e4}{W/m^2}$，在 $2a =
\qty{0.046}{m^2}$ 上，每米 \qty{4}{kW/m}，即每米 $3.5\%$，约 \qty{0.15}{dB/m}。
(d) 损耗比电缆低三倍，功率可达兆瓦级，
而非千瓦级。
\end{solution}

\begin{solution}{exo:b2:guided-waves:8}
(a) $f = c\sqrt2/2L = \qty{2.12}{GHz}$。(b) $\delta = \qty{1.4}{\micro m}$，$R_s = \qty{12}{m\ohm}$；
$W = \varepsilon_0E_0^2V/8 = \qty{1.1e-15}{}E_0^2$；$P = \tfrac12R_s(E_0/377)^2 \times 6L^2 =
\qty{2.5e-9}{}E_0^2$：$Q = \omega W/P \approx 6000$。(c) \qty{360}{kHz}；\qty{0.45}{\micro s}。(d)
$Q \sim 10^{10}$：数十兆伏每米的场只需
千瓦级功率维持，而非吉瓦级。
\end{solution}

\begin{solution}{exo:b2:guided-waves:9}
格点 $(m, n, p)$ 满足 $m^2 + n^2 + p^2 \le (2fL/c)^2$，
占据球体的八分之一，即 $\tfrac18\cdot\tfrac43\pi(2fL/c)^3$，乘以两种偏振：
$8\pi L^3f^3/3c^3$。(b) $8\pi f^2/c^3$。(c) $9 \times 10^{-7}$ 每赫兹（约每兆赫一个，
在 \qty{1}{GHz} 处）；在 \qty{500}{THz} 处为 $2 \times 10^5$ 每赫兹。(d) 当模式
间距超过其宽度时，模式是分立的；在微波炉中，每
隔约 \qty{10}{MHz} 一个，位于 \qty{2.45}{GHz} 附近——在
几个百分点之内有数十个。
\end{solution}

\begin{solution}{exo:b2:guided-waves:10}
(a) 由 $\operatorname{\vect{curl}}\vect E = -\partial_t\vect B$：
\[
\underline B_x = -\frac{k_g}\omega\underline E , \qquad
\underline B_z = \frac{\iu\pi}{a\omega}E_0\cos\frac{\pi x}a\,\eu^{\iu(\omega t - k_gz)} ,
\]
且 $\partial_xB_x + \partial_zB_z = -(k_g\pi/a\omega)E_0\cos(\pi x/a) + (k_g\pi/a\omega)E_0
\cos(\pi x/a) = 0$。(b) $(\operatorname{\vect{curl}}\vect B)_y = \partial_zB_x - \partial_xB_z =
(\iu/\omega)(k_g^2 + \pi^2/a^2)\underline E = (\iu\omega/c^2)\underline E$：$k_g^2 = \omega^2/c^2 - \pi^2/a^2$。
(c) 宽壁：$\vect j_s = \pm(B_z\vect e_x - B_x\vect e_z)/\mu_0$——沿 $z$ 方向的分量
按 $\sin(\pi x/a)$ 变化，横向分量按 $\cos(\pi x/a)$ 变化（在中心处为零）；窄壁：
仅沿 $y$ 方向。(d) 沿中心的纵向缝隙只切断横向电流，
而该处横向电流为零；横向缝隙切断
纵向电流并辐射，即缝隙波导天线。
\end{solution}

\begin{solution}{exo:b2:guided-waves:11}
(a) $2a\sqrt{n_1^2 - n_2^2}/\lambda = 2 \times 50 \times 0.242/1.5 \approx 16$，每种偏振。
(b) $a < \lambda/2\mathrm{NA} = \qty{3.1}{\micro m}$。(c) 数值孔径 $\mathrm{NA} =
0.12$ 的圆形纤芯，在 $d = 2.405\lambda/\pi\mathrm{NA} = \qty{10}{\micro m}$ 以下为单模：
即 \qty{9}{\micro m} 的标准规格。(d) 掠射光线总是被全反射，
因此最低阶模在每个频率下都存在，其倏逝尾部仅
向包层中延伸。
\end{solution}

\begin{solution}{exo:b2:guided-waves:12}
(a) 偏轴光线大部分时间处于折射率较低处，
该处光速较快：较长的路径得到了补偿。(b) $n(r)\cos\theta = n_1\cos\theta_0$，
其中 $n \approx n_1(1 - \Delta r^2/a^2)$ 与 $\cos\theta \approx 1 - \theta^2/2$：$(\dd r/\dd z)^2 =
\theta_0^2 - 2\Delta r^2/a^2$，为简谐振动 $r = r_0\sin(\sqrt{2\Delta}\,z/a)$。
(c) $2\pi a/\sqrt{2\Delta} = \qty{1.1}{mm}$。(d) $n_1L\Delta^2/2c = \qty{0.25}{ns/km}$，相对
\qty{49}{ns/km}，优两百倍。
\end{solution}

\begin{solution}{pb:b2:guided-waves:1}
\textbf{1.} $c/2a = \qty{1.74}{GHz}$；TE$_{20}$ 与 TE$_{01}$：\qty{3.49}{GHz}：仅
TE$_{10}$ 在 \qty{2.45}{GHz} 处传播。

\textbf{2.} $\sqrt{1 - (1.74/2.45)^2} = 0.70$: $\lambda_g = \qty{17.4}{cm}$, $v_\varphi = 1.42c$,
$v_g = 0.70c = \qty{2.1e8}{m/s}$.

\textbf{3.} $E_0^2 = 4P/\varepsilon_0c\,ab\sqrt{\cdot} = 3200/6.9 \times 10^{-6}$: $E_0 = \qty{21}{kV/m}$,
约为击穿场强的百分之一。

\textbf{4.} $j_s \approx E_0/\mu_0c = \qty{57}{A/m}$; $\delta = \qty{10}{\micro m}$, $R_s = \qty{0.1}{\ohm}$;
$\tfrac12R_sj_s^2 \times 2a = \qty{28}{W/m}$：$3.5\%$ 每米——钢材损耗大，
而波导很短。

\textbf{5.} $f = 1.5 \times 10^8\sqrt{11.1(m^2 + n^2) + 25p^2}$，频率在 $\pm3\%$ 以内
接近 \qty{2.45}{GHz}：$(1, 2, 3)$，$(2, 1, 3)$ 在 \qty{2.51}{GHz}，$(4, 0, 2)$，$(0, 4, 2)$，
$(4, 3, 0)$，$(3, 4, 0)$ 在 \qty{2.50}{GHz}——共约六个。

\textbf{6.} 半波长 $a/m$、$b/n$、$d/p$ 的间距为 $7$ 至 \qty{10}{cm}。转盘
带动食物穿过极大与极小区域，搅拌器（一个
旋转的金属扇叶）打乱各模式的分布。

\textbf{7.} $f/Q = \qty{12}{MHz}$：有载共振相互叠合，
连成一片连续谱；磁控管总能找到一个可供馈入的模式。

\textbf{8.} $W = QP/\omega = 3000 \times 800/1.54 \times 10^{10} = \qty{0.16}{J}$; $u = \qty{8.7}{J/m^3}$;
$E_0 \approx \sqrt{4u/\varepsilon_0} = \qty{2}{MV/m}$——接近击穿，产生电弧；反射的
功率返回磁控管，由假负载
或环行器只提供少许保护——这也就是“绝不可空载运行”的原因。

\textbf{9.} $0.2/2.1 \times 10^8 = \qty{1}{ns}$.

\textbf{10.} 峰值 \qty{1.6}{kW}；两脉冲之间，场在
$Q/\omega = \qty{13}{ns}$ 内振荡衰减，谐振腔大部分时间是空的。

\textbf{11.} $\ell \approx a/\pi = \qty{0.48}{mm}$；$\eu^{-1/0.48}$：按振幅计为 \qty{-18}{dB}，
功率则为 \qty{-36}{dB}。

\textbf{12.} $\lambda/4 = \qty{3.1}{cm}$.

\textbf{13.} \qty{5}{mW/cm^2} $\times$ \qty{100}{cm^2} = \qty{0.5}{W}: $0.06\%$.

\textbf{14.} $f_1 \approx c/2d = \qty{3.75}{GHz} > \qty{2.45}{GHz}$：低于截止频率；$\ell =
\qty{1.7}{cm}$, $\eu^{-30/1.7}$: \qty{-150}{dB}.

\textbf{15.} 壁面电流必须连续流入网格；
若网格接合处有间隙，它便是切断电流的缝隙——
一个向外辐射的缝隙天线。

\textbf{16.} $\mathrm{NA} = \sqrt{1.465^2 - 1.460^2} = 0.12$，\qty{7}{\degree}：很窄的
锥形光束，需要激光器和透镜。

\textbf{17.} $(n_1L/c)(n_1/n_2 - 1) = \qty{17}{ns/km}$；\qty{170}{ns} 在 \qty{10}{km} 上：
\qty{6}{Mbit/s}.

\textbf{18.} $\tfrac12(\pi d\,\mathrm{NA}/\lambda)^2 = 75$；对于 \qty{9}{\micro m}：$2.4$——即
单模区（$V = \pi d\,\mathrm{NA}/\lambda = 2.2 < 2.4$）。

\textbf{19.} $\Delta x_0 = (c/n)\tau = \qty{1}{cm}$: $t_{\text{展宽}} = 10^{-4}/0.2 = \qty{0.5}{ms}$,
\qty{100}{km}：补偿前，该距离上可达 \qty{10}{Gbit/s}。

\textbf{20.} \qty{20}{dB}，即 $100$ 倍：\qty{10}{\micro W}；$6000/80 = 75$
个放大器。

\textbf{21.} $f = \qty{1.9e14}{Hz}$，$hf = \qty{0.8}{eV}$；每秒 $8 \times 10^{15}$ 个
光子；输入端每比特 $8 \times 10^5$ 个，经 \qty{100}{km} 后为 $8 \times 10^3$ 个。

\textbf{22.} 两种损耗之和在 \qty{1.55}{\micro m} 附近最小
（\qty{0.2}{dB/km}）。

\textbf{23.} $R = ((1.465 - 1)/2.465)^2 = 3.6\%$：每个端面 \qty{0.16}{dB}，两次
出现在连接器处，另有间隙的干涉；凝胶或物理接触
可排除空气。

\textbf{24.} 弯道平缓时，光线仍以超过
临界角的入射角射到边界上；急弯时，射到外壁的入射角
便低于临界角，光就漏出去。

\textbf{25.} 烤箱：金属壁面反射，$r = -1$；带宽由
谐振腔的模式决定；损耗在钢壁趋肤层中。光纤：全
反射发生在玻璃–玻璃界面上；带宽由模式与
色度色散决定；损耗来自散射与吸收，\qty{0.2}{dB/km}。
\end{solution}
